\documentclass[a4paper,11pt]{article}
\usepackage[T1]{fontenc}
\usepackage[utf8]{inputenc}
\usepackage[italian]{babel}
\usepackage[top=2.5cm,bottom=2.5cm,left=2.5cm,right=2.5cm]{geometry}
\usepackage{hyperref}
\usepackage{xcolor}
\definecolor{accent}{HTML}{1F4E78}
\hypersetup{colorlinks=true, urlcolor=accent}
\pagestyle{empty}
\setlength{\parindent}{0pt}
\setlength{\parskip}{10pt}

\begin{document}

Gentili Signore e Signori,

vi scrivo per candidarmi alla posizione di Software Developer Angular \& Analyst presso Finomnia, per la sede di Brescia. Ho tre anni di esperienza come sviluppatore full-stack in diverse aziende europee, con particolare attenzione allo sviluppo frontend, alle integrazioni API e all'analisi dei requisiti funzionali, e sono motivato a portare questa esperienza nel team Digital Solutions \& Front-End Applications.

Per FlowSpace, marketplace londinese, ho sviluppato il frontend in TypeScript e React/Next.js e ho automatizzato la migrazione dai sistemi legacy alla nuova piattaforma, mappando la struttura esistente sulla nuova architettura e migrando i dati entità per entità tramite integrazioni REST e file CSV, con il prodotto sempre attivo durante il passaggio; la piattaforma è cresciuta da 2.000 a oltre 4.500 utenti.

Per Clubee, in Lussemburgo, ho integrato le API della piattaforma su tutto il frontend e mi sono occupato di unit testing e controllo qualità delle nuove funzionalità, un percorso vicino alle attività di collaudo e avviamento delle soluzioni che descrivete nell'annuncio.

La mia esperienza frontend più recente è in TypeScript con React e Next.js, non ancora in Angular; conosco bene i concetti di componenti, servizi e binding tipizzato e mi aspetto un transfer rapido verso Angular 2+, che sono pronto ad affrontare fin da subito con impegno nello studio.

Sono cittadino indiano con residenza di lungo periodo UE in Italia, quindi autorizzato a lavorare in Italia senza necessità di sponsorizzazione. Sono disponibile a trasferirmi a Brescia e a iniziare da subito. Il mio livello di italiano è intermedio, sufficiente per lavorare e comunicare in team, mentre l'inglese è il mio livello professionale principale.

Resto a disposizione per un colloquio.

Cordiali saluti,\\
Prabhavit Gupta

\end{document}
